\chapter{从卵细胞到出生}\label{ch:g8:from-egg-to-birth}

\Cref{ch:g5:human-reproduction}以故事形式讲述了九个月；
现在你已掌握讲好这一故事所需的全部工具，
并能熟练运用：\omterm{def:g5:first-look-at-cells:cell}{细胞}及其分裂、两种\omterm{def:g8:sexual-reproduction:gametes}{配子}及其相遇之处、
\omterm{def:g8:puberty:hormones}{激素}，以及\omterm{prop:g4:journey-of-food:absorb}{血液}的输送。本章将借助这些机制，
重新讲述人类的发育过程，让其机制清晰可见——
从一个\omterm{def:g5:first-look-at-cells:cell}{细胞}到三万亿个\omterm{def:g5:first-look-at-cells:cell}{细胞}，历时四十周。

\section{第一周：从受精卵到“寄居者”}

\begin{proposition}[沿输卵管下行]\label{prop:g8:from-egg-to-birth:firstweek}
\omterm{def:g5:flower-to-fruit:steps}{受精}发生在\omterm{def:g8:reproductive-systems:female}{输卵管}中
（\cref{def:g8:reproductive-systems:female}）同时启动
两个“时钟”：
\begin{enumerate}
  \item \emph{分裂}：\omterm{prop:g8:sexual-reproduction:core}{受精卵}在一天之内开始分裂，
        此后不断重复分裂
        （\cref{ex:g5:first-look-at-cells:one}）——两个、
        四个、八个\omterm{def:g5:first-look-at-cells:cell}{细胞}，到第四天形成数十个\omterm{def:g5:first-look-at-cells:cell}{细胞}的桑葚胚，
        一边沿\omterm{def:g8:reproductive-systems:female}{输卵管}下行，
        向\omterm{prop:g5:human-reproduction:plan}{子宫}方向移动；
  \item \emph{着床}：约在第七天，这个小球
        到达\omterm{prop:g5:human-reproduction:plan}{子宫}，并嵌入
        已准备好的\omterm{prop:g5:human-reproduction:plan}{子宫}内膜
        （\cref{prop:g8:reproductive-systems:cycle}中的事件
        4）——“寄居者”就位，循环停止，
        \omterm{def:g8:puberty:hormones}{激素}向整个系统发出信号：\omterm{def:g5:human-reproduction:pregnancy}{怀孕}开始。
\end{enumerate}
\end{proposition}
\admitted

\begin{example}[早期分裂靠什么维持]\label{ex:g8:from-egg-to-birth:provisions}
在第一周里，这位“旅行者”不进餐，
也不接受任何供给：它完全依靠\omterm{def:g4:animal-reproduction:sexual}{卵细胞}自带的
养料——这正是
\cref{ex:g8:reproductive-systems:egg}中储备养料的用意。
\omterm{def:g8:reproductive-systems:male}{精子}只带来了父方的\omterm{def:g6:cell-unit-of-life:parts}{细胞核}，几乎别无他物；
\omterm{def:g4:animal-reproduction:sexual}{卵细胞}则带来了另一个\omterm{def:g6:cell-unit-of-life:parts}{细胞核}，\emph{还}带来了养料库。
着床的意义之一，在于养料库
渐趋耗尽，而“寄居者”改由母体“客房送餐”。
\end{example}

\section{胚胎、胎盘、胎儿}

\begin{definition}[胎盘]\label{def:g8:from-egg-to-birth:placenta}
着床后的\omterm{def:g5:first-look-at-cells:cell}{细胞}团同时发育出两种结构：
一是\emph{\omterm{def:g5:human-reproduction:pregnancy}{胚胎}}本身，另一是其
生命支持界面——即\emph{胎盘}\index{胎盘}，呈盘状
长入\omterm{prop:g5:human-reproduction:plan}{子宫}壁内，母体的\omterm{prop:g4:journey-of-food:absorb}{血液}与
\omterm{def:g5:human-reproduction:pregnancy}{胚胎}的\omterm{prop:g4:journey-of-food:absorb}{血液}紧密相邻，却\emph{互不
混合}。其广阔而薄的表面——依
\cref{prop:g7:how-animals-breathe:spec}所述规格
再次搭建，\omterm{def:g7:digestion-and-nutrients:families}{营养物质}与\omterm{def:g7:what-is-respiration:respiration}{氧气}经此转运至\omterm{def:g5:human-reproduction:pregnancy}{胚胎}，
废物与\omterm{def:g7:what-is-respiration:respiration}{二氧化碳}则反向穿过
（\cref{prop:g5:human-reproduction:cord}中的\omterm{prop:g5:human-reproduction:cord}{脐带}
连接胎盘与\omterm{def:g5:human-reproduction:pregnancy}{胚胎}）。胎盘对“乘客”而言，兼具肺
、肠与肾的功能——三道“边界”
全都集于一个器官之中。
\end{definition}

\begin{proposition}[借助工具重新审视时间表]\label{prop:g8:from-egg-to-birth:timetable}
\Cref{ex:g5:human-reproduction:timetable}中的各个里程碑，
现逐一加以解释：
\begin{itemize}
  \item 第 1--8 周，\emph{\omterm{def:g5:human-reproduction:pregnancy}{胚胎}}期：身体的基本构造
        已经奠定——\omterm{def:g5:first-look-at-cells:cell}{细胞}分裂、迁移并
        \emph{分化}
        （\cref{rem:g6:cell-unit-of-life:twoways}）
        形成\omterm{def:g4:heart-and-blood:heart}{心脏}、\omterm{def:g5:brain-in-command:system}{神经}、\omterm{def:g3:bones-and-muscles:skeleton}{骨骼}等器官及其余部分；\omterm{def:g4:heart-and-blood:heart}{心脏}自
        第四周起开始搏动，因为这套“养料与\omterm{def:g8:from-egg-to-birth:placenta}{胎盘}”的供应体系
        需要尽早拥有它的“泵”；
  \item 第 9--40 周，\emph{\omterm{def:g5:human-reproduction:pregnancy}{胎儿}}期：构造已成，
        生长与成熟——\omterm{def:g5:brain-in-command:system}{大脑}自行连接，
        肺则在无空气时预先“演练”，
        （\cref{exo:g5:human-reproduction:11}中已解释
        其首次啼哭的“首演”），第四个月的胎动
        表明\omterm{def:g3:bones-and-muscles:muscle}{肌肉}开始工作；
  \item 贯穿始终：\omterm{def:g8:puberty:hormones}{激素}使\omterm{prop:g5:human-reproduction:plan}{子宫}保持平静，
        循环暂停——此即
        \cref{def:g8:puberty:hormones}所述的化学机制，
        统筹孕期的运作。
\end{itemize}
\end{proposition}
\admitted

\begin{example}[为何孕早期最需谨慎]\label{ex:g8:from-egg-to-birth:earlycare}
\Cref{prop:g5:human-reproduction:care}曾警告说，
烟草与酒精会传递给\omterm{def:g5:human-reproduction:pregnancy}{胎儿}；而时间表
则指明最关键的时段：在奠定基本构造的\omterm{def:g5:human-reproduction:pregnancy}{胚胎}期，
即每个器官还只是一小群正在分化的\omterm{def:g5:first-look-at-cells:cell}{细胞}时，
一条被扰乱的指令便可能在整个构建过程中引发
连锁“回声”——而这恰恰是\omterm{def:g5:human-reproduction:pregnancy}{怀孕}
可能在无人察觉中度过的几周。因此医生的
算术很直白：护理不应始于首次超声检查，
而应始于可能\omterm{def:g5:human-reproduction:pregnancy}{怀孕}的那一刻。
\end{example}

\section{出生与转换}

\begin{proposition}[出生：由激素调控的转换]\label{prop:g8:from-egg-to-birth:birth}
约在第四十周，\omterm{def:g8:puberty:hormones}{激素}信号——\omterm{def:g5:human-reproduction:pregnancy}{胎儿}与母体的
“对话”——结束了\omterm{prop:g5:human-reproduction:plan}{子宫}长久的平静：
\begin{enumerate}
  \item \emph{分娩}：\omterm{prop:g5:human-reproduction:plan}{子宫}是人体最强壮的
        \omterm{def:g3:bones-and-muscles:muscle}{肌肉}（\cref{rem:g8:reproductive-systems:living}
        曾提及其较温和的工作），以渐强的
        收缩波打开通道，将\omterm{def:g5:human-reproduction:pregnancy}{胎儿}推向
        产道；
  \item \emph{出生}：通常头先娩出
        （\cref{ex:g5:human-reproduction:birth}）；
  \item \emph{转换}：第一声啼哭充满了
        从未使用过的肺；\omterm{prop:g5:human-reproduction:cord}{脐带}的血流终止，它
        被剪断；数分钟之内，新生儿便依靠自身的
        肺、即将启用的肠道与\omterm{def:g7:removing-wastes:kidneys}{肾脏}——\omterm{def:g8:from-egg-to-birth:placenta}{胎盘}
        此前覆盖的每一道“边界”，如今都由它自己承担。
        （\cref{exo:g5:human-reproduction:11}，
        现以完整图景给出）；
  \item \emph{\omterm{def:g8:from-egg-to-birth:placenta}{胎盘}娩出}：\omterm{def:g8:from-egg-to-birth:placenta}{胎盘}完成使命后
        脱离并随之排出——这是人体唯一
        为一次性使用而构建、用后即弃的器官。
\end{enumerate}
\end{proposition}
\admitted

\begin{example}[新生儿的第一次“体检”]\label{ex:g8:from-egg-to-birth:audit}
助产士在婴儿出生后数分钟内为其评分——
\omterm{def:g7:what-is-respiration:respiration}{呼吸}、心跳、肤色、活动、哭声。用本书的眼光读这张
清单：它正是
\cref{prop:g7:removing-wastes:complete}中的
供给与排出循环，在其首次独立运行时接受检验。
每个人最先通过的考试，是一场生理学测试。
\end{example}

\begin{method}[解释妊娠的任一阶段]\label{met:g8:from-egg-to-birth:explain}
若被问及\omterm{def:g5:human-reproduction:pregnancy}{怀孕}九个月中的任一时刻，可用
四个工具作答：
\begin{enumerate}
  \item \emph{\omterm{def:g5:first-look-at-cells:cell}{细胞}}：分裂与分化
        在做什么？
  \item \emph{供给}：是“储藏室”，还是\omterm{def:g8:from-egg-to-birth:placenta}{胎盘}——什么
        可以通过，什么不可以？
  \item \emph{\omterm{def:g8:puberty:hormones}{激素}}：正在被维持，还是正在被下达指令？
  \item \emph{照护}：这一阶段对母亲的习惯与
        医生的日程有何要求？
\end{enumerate}
\end{method}

\section{练习}

\begin{exercise}[$\star$]\label{exo:g8:from-egg-to-birth:1}
为第一周安排时间线：分裂的开始、桑葚
胚、着床。
\end{exercise}

\begin{exercise}[$\star$]\label{exo:g8:from-egg-to-birth:2}
第一周的“旅行者”以什么为养料？什么
结束了这种安排？
\end{exercise}

\begin{exercise}[$\star$]\label{exo:g8:from-egg-to-birth:3}
什么是\omterm{def:g8:from-egg-to-birth:placenta}{胎盘}？它双向各传递什么物质，
又有什么物质从不通过？
\end{exercise}

\begin{exercise}[$\star$]\label{exo:g8:from-egg-to-birth:4}
将四十周划分为\omterm{def:g5:human-reproduction:pregnancy}{胚胎}期与\omterm{def:g5:human-reproduction:pregnancy}{胎儿}期，并说明每
一时期的任务。
\end{exercise}

\begin{exercise}[$\star$]\label{exo:g8:from-egg-to-birth:5}
为什么\omterm{def:g4:heart-and-blood:heart}{心脏}在第四周就开始搏动，且早于
所有器官？
\end{exercise}

\begin{exercise}[$\star$]\label{exo:g8:from-egg-to-birth:6}
按顺序列出分娩的四个编号事件。
\end{exercise}

\begin{exercise}[$\star\star$]\label{exo:g8:from-egg-to-birth:7}
\omterm{def:g8:from-egg-to-birth:placenta}{胎盘}承担了哪三个器官的功能？指出
各自在前文中对应的章节。
\end{exercise}

\begin{exercise}[$\star\star$]\label{exo:g8:from-egg-to-birth:8}
为什么\omterm{def:g5:human-reproduction:pregnancy}{胚胎}期最需要照护？请依据
\cref{ex:g8:from-egg-to-birth:earlycare}给出两个理由。
\end{exercise}

\begin{exercise}[$\star\star$]\label{exo:g8:from-egg-to-birth:9}
将新生儿的评分表按
\cref{prop:g7:removing-wastes:complete}中的循环来读：
将\omterm{def:g7:what-is-respiration:respiration}{呼吸}、心跳和肤色分别与其对应的系统相匹配。
\end{exercise}

\begin{exercise}[$\star\star$]\label{exo:g8:from-egg-to-birth:10}
重访同卵双胞胎
（\cref{rem:g5:human-reproduction:twins}）：在
\cref{prop:g8:from-egg-to-birth:firstweek}的哪个时间点
发生单个\omterm{prop:g8:sexual-reproduction:core}{受精卵}的分裂？时间安排为何使得
他们“如同复制品般相似”？
\end{exercise}

\begin{exercise}[$\star\star$]\label{exo:g8:from-egg-to-birth:11}
以\cref{met:g8:from-egg-to-birth:explain}分析第
四个月——即胎动初现的月份。
\end{exercise}

\begin{exercise}[$\star\star\star$]\label{exo:g8:from-egg-to-birth:12}
“出生是边界的转换。”请论证：列出\omterm{def:g8:from-egg-to-birth:placenta}{胎盘}
处理过的每一次交换，以及在第一声
啼哭时接替其工作的器官——并说明这一
转换以分钟为单位所体现的显著时效性。
\end{exercise}

\section{问题：助产士的时间线}

\begin{problem}[{周末问题——逐次
产检的一次妊娠}]\label{pb:g8:from-egg-to-birth:1}
助产士为一次（虚构的）妊娠建立的档案：
第 8 周建档，第 12 周和第 20 周超声检查，第 28
周和第 36 周常规检查，第 40 周分娩。请用本章的
工具阅读这份档案。

\medskip\noindent\textbf{第一部分——档案的
早期记录。}
\begin{enumerate}
  \item 妊娠通过\omterm{def:g8:puberty:hormones}{激素}检测得到确认，
        早于首次就诊数周。
        这种检测读取的是哪一种信号？它从
        \cref{prop:g8:from-egg-to-birth:firstweek}
        中的哪一事件起开始存在？
  \item 第 8 周时，助产士记录“\omterm{def:g5:human-reproduction:pregnancy}{胚胎}基本完成，
        大小约如一粒豆子”。请用时间线的术语翻译此句——
        此处结束的是哪一
        个时期？
  \item 建档检查的建议首先是：
        不饮酒、不吸烟，并补充叶酸，即“建造者的维生素”，
        “最好从孕前开始”。
        论证这三条建议，请依据
        \cref{ex:g8:from-egg-to-birth:earlycare}。
  \item 第 12 周的超声检查显示了
        从第 4 周开始的心跳。从供给的角度看，为何这么早？
\end{enumerate}

\medskip\noindent\textbf{第二部分——孕中期。}
\begin{enumerate}[resume]
  \item 第 20 周的超声检查测量了每个器官以及
        \omterm{def:g8:from-egg-to-birth:placenta}{胎盘}的位置。请说明\omterm{def:g5:human-reproduction:pregnancy}{胎儿}期的
        任务，以及正在接受检查的\omterm{def:g8:from-egg-to-birth:placenta}{胎盘}
        三项职能。
  \item 母亲报告约在第 18 周时出现首次胎动。
        此时正在显现的是哪一套机制
        （\cref{prop:g5:brain-in-command:loop}
        给出了该链条）？
  \item 第 28 周时，助产士检查母亲的
        血糖与血压。为何母亲的循环
        承受双重负荷？它正穿越的
        是谁的边界？
  \item 第 30--36 周：\omterm{def:g5:human-reproduction:pregnancy}{胎儿}主要在生长，肺在
        进行演练。演练的是什么，
        又是为了哪一个单一时刻？
\end{enumerate}

\medskip\noindent\textbf{第三部分——最后一页。}
\begin{enumerate}[resume]
  \item 第 40 周：分娩。说出该\omterm{def:g3:bones-and-muscles:muscle}{肌肉}的名称、
        它的两项任务，以及下达这项工作的
        是谁（\cref{prop:g8:from-egg-to-birth:birth}）。
  \item 档案中的分娩记录为：“出生后即刻有力啼哭；
        \omterm{prop:g5:human-reproduction:cord}{脐带}夹闭并剪断；\omterm{def:g8:from-egg-to-birth:placenta}{胎盘}完整。”
        请结合这一转变逐条核查这三项记录。
  \item 新生儿在按分钟记分的评分表中得分很高。
        请说明该检查表证明了什么，
        按\cref{ex:g8:from-egg-to-birth:audit}
        的术语作答。
  \item 请用一句话为父母总结这份档案：四十周
        由哪一个\omterm{def:g6:cell-unit-of-life:unicellular}{单细胞}出发，借助何种机制，
        建成了什么。
\end{enumerate}
\end{problem}
